% !TEX root = ../main.tex


\section{Introduction}

Federated learning~\cite{yang2019federated} facilitates collaborations among a set of clients and preserves their privacy so that the clients can achieve better machine learning performance than individually working alone. The underlying idea is to collectively learn from data from all clients. The initial idea of federated learning starts from aggregating models from clients to achieve a global model so that the global model can be more general and capable. The effectiveness of this global collaboration theme that is not differentiating among all clients highly depends on the data distribution among clients.  It works well on IID data, that is, clients are similar to each other in their private data distribution.

\nop{
With the advances in internet of things (IoT) and 5G technology, 
}
\nop{
Huge amount of edge devices (i.e., clients), such as smart phones, wearable devices and autonomous vehicles, are generating big private data in a surprisingly fast speed~\cite{poushter2016smartphone}.
Due to the ever-growing concerns and restrictions on data privacy,
many federated learning frameworks have been proposed to collaboratively conduct machine learning tasks for all clients while keeping their data privacy intact~\cite{mcmahan2016communication, li2018federated, zhao2018federated, erlingsson2014rappor, hard2018federated, yang2018applied}.
}

In many application scenarios where collaborations among clients are needed to train machine learning models, data are unfortunately not IID.  For example, consider the cases of personalized cross-silo federated learning~\cite{kairouz2019advances}, where there are tens or hundreds of clients and the private data of clients may be different in size, class distributions and even the distribution of each class.  Global collaboration without considering individual private data often cannot achieve good performance for individual clients. 

\nop{Some federated learning methods, such as FedAvg and FedProx, try to fix the problem by conducting an additional fine-tuning step after a global model is trained~\cite{ben2010theory, cortes2014domain, mansour2020three, mansour2009domain,  schneider2019mass, wang2019federated}. 
While those methods work in some cases, they cannot solve the problem systematically as demonstrated in our experimental results (e.g., dataset CIFAR100 in Table~\ref{tab:result}).}

Some federated learning methods try to fix the problem by conducting an additional fine-tuning step after a global model is trained~\cite{ben2010theory, cortes2014domain, mansour2020three, mansour2009domain,  schneider2019mass, wang2019federated}.While those methods work in some cases, they cannot solve the problem systematically as demonstrated in our experimental results (e.g., data set CIFAR100 in Table~\ref{tab:result}).

We argue that the fundamental bottleneck in personalized cross-silo federated learning with non-IID data is the misassumption of one global model can fit all clients. Consider the scenario where each client tries to train a model on customers' sentiments on food in a country.  Different clients collect data in different countries. Obviously, customers' reviews on food are likely to be related to their cultures, life-styles, and environments. Unlikely there exists a global model universally fitting all countries.  Instead, pairwise collaborations among countries that share similarity in culture, life-styles, environments and other factors may be the key to accomplish good performance in personalized cross-silo federated learning with non-IID data.

\nop{
 has attracted intensive attention from large companies and institutes who want to participate as clients in the learning framework to collaboratively train strong personalized machine learning models without exposing their private data~\cite{kairouz2019advances}.
For example, multiple banks can collaboratively train robust machine learning models to accurately evaluate the credit scores of users~\cite{yang2019federated}, 
different insurance companies may collaboratively train strong fraud detection models to identify fraudsters more accurately~\cite{li2019federated}, and different hospitals can build better diagnostic models together without exposing sensitive private information of patients~\cite{brisimi2018federated, courtiol2019deep}.
}
\nop{
One of the biggest challenges for personalized cross-silo federated learning is that the data is non-IID across different clients~\cite{kairouz2019advances, li2018federated, zhao2018federated}.
Since it is difficult for a single model to properly fit the non-IID data of all clients~\cite{kairouz2019advances}, existing global federated learning methods~\cite{mcmahan2016communication, li2018federated, ji2019learning, yurochkin2019bayesian, wang2020federated} that train a single global model cannot achieve a good personalized performance on every client.
}

\nop{
As a result, it is more reasonable to conduct \textit{personalized federated learning}, such that every client can train a strong personalized model of its own by effectively collaborating with the other clients without exposing its private data.
}

\nop{
Since the non-IID data of two clients ca either be similar or different, an effective inter-client collaboration should allow closely related clients with similar data distributions to have stronger collaboration than those with different data distributions~\cite{smith2017federated}.


The performance of personalized cross-silo federated learning largely depends on the effectiveness of inter-client collaboration, which, however, is a challenging goal to achieve because the inter-client collaboration relationship varies a lot for different pairs of clients
when the private data of clients is non-IID.
As illustrated later in Section~\ref{sec:relatedworks}, 
existing global federated learning methods~\cite{ji2019learning,  li2018federated,
mcmahan2016communication, wang2020federated,yurochkin2019bayesian} that train a single global model cannot achieve a good personalized performance on every client, because it is difficult for a single model to properly fit the non-IID data of all clients~\cite{kairouz2019advances}.
Most existing personalized federated learning methods~\cite{deng2020adaptive, fallah2020personalized,  hanzely2020federated, karimireddy2019scaffold, mansour2020three, schneider2019mass, wang2019federated} use a single global model to conduct a universal collaboration between all clients. These methods can distinguish the amount of contribution made by each client to the global model, but they cannot describe the underlying pair-wise collaboration relationships between different pairs of clients. 
This reduces the effectiveness of inter-client collaboration for clients holding non-IID private data, which further limits the performance of personalized federated learning.
}

\nop{
cannot effectively identify the underlying collaboration relationships between clients.

Most of the local customization methods use a single global model to conduct a universal collaboration involving all clients. This universal collaboration can only distinguish the amount of contribution made by each client, but cannot describe the pair-wise collaboration relationships between clients.
As a result, the local customization methods mentioned above cannot effectively discover the underlying collaboration relationships between different pairs of clients, which largely limits the effectiveness of inter-client collaboration when the data is non-IID across different clients. 
}

\nop{
As illustrated later in Section~\ref{sec:relatedworks}, existing methods~\cite{mansour2020three, schneider2019mass, wang2019federated, deng2020adaptive, hanzely2020federated, fallah2020personalized, nichol2018first, chen2018federated} train personalized models for different clients by local customization, that is, to let the client customize the collaboratively learnt global model using its private local data.
}

\nop{
customize the global model for a client using its local data.
However, without knowing the private data distributions of clients, these methods cannot effectively identify the underlying collaboration relationships between clients, which reduces the effectiveness of collaboration and largely limits their performance.

As a result, these methods can only 

most of these methods collaboratively produce a global model


that take simple average of client models cannot identify the underlying collaboration relationships between clients.
This reduces the effectiveness of collaboration and largely limits their performance.






The performance of federated learning is largely determined by 

The effectiveness of inter-client collaboration largely determines the performance of federated learning. 





Since a single model cannot properly fits the non-IID data of all clients~\cite{kairouz2019advances}, existing federated learning frameworks that focus on training a single global model for all clients cannot achieve a good personalized performance on each of the clients.
}

\nop{Classical centralized learning methods are only applicable when the data from all clients are put together in a central server. 
However, due to the ever-growing concerns and restrictions on data privacy, it is very difficult to gather the private data of all clients~\cite{yang2019federated}.
As a result, federated learning rises as an attractive framework to collaboratively conduct machine learning tasks for all clients while keeping their data privacy intact~\cite{mcmahan2016communication, li2018federated, zhao2018federated}.
}
\nop{
federated learning has emerged as an attractive distributed learning framework to unleash the true power of the huge private data scattered in various edge devices, such as smart phones, smart watches, and autonomous vehicles.

A large number of edge devices gives demands for federated learning.

A general definition of federated learning.
Federated learning (FL) is a machine learning setting where many clients (e.g. mobile devices or whole organizations)
collaboratively train a model under the orchestration of a central server (e.g. service provider),
while keeping the training data decentralized.


Federated learning has many practical applications, such as google, ios 13 and so on.
}

\nop{
Many federated learning frameworks have been proposed to tackle privacy-preserving machine learning problems in a broad range of applications~\cite{erlingsson2014rappor, hard2018federated, yang2018applied}. 

However, the practical performance of most existing methods are limited due to the following challenges~\cite{mcmahan2016communication, konevcny2015federated, li2019federated, konevcny2016federated}.
}

\nop{
many federated learning frameworks has been  federated learning often faces the following practical challenges.
}

\nop{
the data distribution is non-IID across different clients. 
}

\nop{
\textbf{Non-IID data}: 
In contrast to the independent and identically distributed (IID) data in centralized training,
the data in federated learning is usually non-IID across different clients~\cite{li2018federated, zhao2018federated, kairouz2019advances}.
Many existing federated learning frameworks focus on training a global model to achieve a good performance on the aggregated data of all clients.
However, when the data is non-IID across different clients, it is difficult to train a single global model that properly fits the data of all clients~\cite{kairouz2019advances}. Instead, it is more reasonable to conduct \textit{personalized federated learning} that trains a personalized model for each client~\cite{smith2017federated}.
}

\nop{
This global performance, however, is hard to be good when the data distributions are non-IID across different clients, because 
}

\nop{
has no relevance if the global
model is expected to be subsequently personalized before being
put to use.
Personalized models usually show better performance for
individual clients than global or local models. 
}

\nop{Take google keyboard users as an example~\cite{hard2018federated, yang2018applied}, users with similar user habits generate data with similar distributions. These closely related users can get better performances if they collaborate more.}

\nop{
\textbf{Unknown collaboration relationship}: 
The effectiveness of inter-client collaboration largely determines the performance of federated learning.
Intuitively, an effective collaboration should allow closely related clients with similar data distributions to have stronger collaborations than those with different data distributions.
However, without knowing the private data distributions of clients, classical federated learning methods that take simple average of client models cannot identify the underlying collaboration relationships between clients.
This reduces the effectiveness of collaboration and largely limits their performance.
}

\nop{
These users form a coherent user-group and should have strong inner-group collaborations with each other.
}

\nop{
Moreover, since different users can have similar user habits, clients can form clusters, such that the clients within the same cluster have similar data distributions.
}

\nop{
the collaboration strength between clients should not always be the same.
Instead, 
}

\nop{
Since different users can have similar user habits, clients can form groups, such that the clients within the same group have similar data distributions.
Take the input method of smart phones as an example, different users type in different ways, but users with similar typing habits can form a user-groups, such that users in the same group generates data with similar distributions.
}

\nop{
With a target of collaboratively training models on edge devices,
federated learning is bound to
}

\nop{
Another practical challenge for personalized federated learning is the \textbf{unreliable operating environment}.
Unlike the stable computing environment on central servers, personalized federated learning is bound to operate in an unreliable environment that involves limited communication bandwidth, unstable network connections, unreliable battery supplies, and highly unbalanced computational resource of clients~\cite{yang2019federated, singh2019detailed}. 
As a result, clients may occasionally drop off-line due to dead batteries or signal loss~\cite{smith2017federated}, and this has been an open challenge for most federated learning frameworks~\cite{mcmahan2016communication, li2018federated, zhao2018federated}.
}

\nop{
some clients may drop off-line due to battery drainage or signal disconnection, 


There can be client drops and stragglers due to limited communication bandwidth.
}

\nop{
More often than not, clients may form groups, where 




the clients' data in federated learning is generated by various end-user devices, thus 
}



\nop{
federated learning uses local data from , leading to many varieties of non-IID data.
}
\nop{
\textbf{mention different distributoin across different clients, and groups of clients with similar distributions.}
For example, ...
With a few exceptions, most prior work is focused on measuring the
performance of the global model on aggregated data instead
of measuring its performance as seen by individual clients.
Global performance, however, has no relevance if the global
model is expected to be subsequently personalized before being
put to use.
Personalized models usually show better performance for
individual clients than global or local models. 

Second, limited communication bandwidth causes client drops and stragllers.
}

\nop{
Propose challenges: non-IID data distribution, unstable communication that causes client drops and stragglers, 
}

\nop{
we tackle the above challenges by formulating the task of personalized federated learning as a multi-task learning problem. The key idea is to allow each client to have its own personalized model, and encourage clients to adaptively collaborate in clusters by forcing similar client models to become more similar during training. 
}

\nop{
allows each client to train a good personalized model by maintaining strong collaborations with its closely related clients.
}
\nop{ 
encourages closely related clients to collaborate more with each other by attentively passing messages (i.e., model parameters) between similar personalized cloud models and local models. by forcing similar client models to become more similar during training. } 

\nop{
by encouraging non-linearly stronger collaborations between the clients with more similar model parameters.
}

\nop{
aggregating the local models of all clients into different cloud models. 
To encourage collaborations between closely related clients,
{FedAMP} iteratively passes non-linearly stronger model aggregation messages between the local models and cloud models that have more similar model parameters.

It conducts inter-client collaboration by attentively passing aggregation messages (i.e., model parameters) between client models and personalized cloud models with similar model parameters in a non-linear fashion.
The non-linear attentive message passing }

\nop{
Instead of using a single global model on the cloud server, our method maintains a personalized cloud model for each client on the cloud server.


FedAMP adopts an attention-inducing function to bind a \textit{personalized cloud model} on the cloud server with each client's local personalized model. In this way, the attentive  can be easily conducted by iteratively passing strong model aggregation messages between similar local models and personalized cloud models.
}

\nop{
In this paper, 
we tackle the challenging personalized federated learning problem by an \textit{attentive message passing mechanism} that adaptively discovers the underlying pair-wise collaboration relationships between clients by iteratively encouraging similar clients to collaborate more.
We make the following contributions.
\textit{First}, we propose a novel method named FedAMP to realize the attentive message passing mechanism.
FedAMP allows each client to own a local personalized model, but does not use a single global model on the cloud server to conduct universal collaboration.
Instead, it maintains a personalized cloud model on the cloud server for each client, and realizes the attentive message passing mechanism by passing strong model-aggregation messages from the personalized model of each client to similar personalized cloud models owned by the other clients.
This adaptively discovers the underlying pair-wise collaboration relationships between clients, and significantly boosts the effectiveness of collaboration.
\textit{Second}, we prove the convergence of FedAMP for both convex and non-convex personalized models.
\textit{Third}, we propose a heuristic method to further improve the practical performance of FedAMP on clients using deep neural networks as personalized models.
\textit{Last}, we conduct extensive experiments to demonstrate the superior performance of the proposed methods.
}

Carrying the above insight, in this paper, we tackle the challenging personalized cross-silo federated learning problem by a novel \textit{attentive message passing mechanism} that adaptively facilitates the underlying pairwise collaborations between clients by iteratively encouraging similar clients to collaborate more.
We make several technical contributions. 

We propose a novel method \emph{federated attentive message passing} (FedAMP) whose central idea is the attentive message passing mechanism.
FedAMP allows each client to own a local personalized model, but does not use a single global model on the cloud server to conduct collaborations.
Instead, it maintains a personalized cloud model on the cloud server for each client, and realizes the attentive message passing mechanism by attentively passing the personalized model of each client as a message to the personalized cloud models with similar model parameters.  Moreover, FedAMP updates the personalized cloud model of each client by a weighted convex combination of all the messages it receives. 
This adaptively facilitates the underlying pairwise collaborations between clients and significantly improves the effectiveness of collaboration.

We prove the convergence of FedAMP for both convex and non-convex personalized models. Furthermore, we propose a heuristic method to further improve the performance of FedAMP on clients using deep neural networks as personalized models. We conduct extensive experiments to demonstrate the superior performance of the proposed methods.

\nop{
, where the weight of each message is positively correlated with the similarity between the message and the personalized cloud model in a non-linear manner
}


\nop{
, 


 to avoid the scaling problem of high dimensional model parameters of clients by measuring their model similarity with cosine similarity
}

\nop{
in challenging federated learning settings, such as non-IID data, dropped clients and dirty data.
}



\nop{
further improve the practical performance of {FedAMP} by measuring the similarity between the model parameters of clients using cosine similarity, which .
}

\nop{
as demonstrated by extensive experiments, the heuristic method further boosts the performance of FedAMP.
}

\nop{
\textbf{Motivation}
\begin{itemize}
    \item Why are we interested in horizontal federated learning?
    \begin{itemize}
        \item Protect data privacy.
        \item Boost performance by collaboration.
        \item Some solid examples of applications of horizontal federated learning.
    \end{itemize}
    
With a few exceptions, most prior work is focused on measuring the
performance of the global model on aggregated data instead
of measuring its performance as seen by individual clients.
Global performance, however, has no relevance if the global
model is expected to be subsequently personalized before being
put to use.
Personalized models usually show better performance for
individual clients than global or local models. 


\nop{
In some cases,
however, personalized models fail to reach the same performance
as local models, especially when differential privacy
and robust aggregation is implemented [7].
}

    \item Why is it emergent to tackle the non-IID problem?
    \begin{itemize}
        \item Non-IID data distribution is the most common problem for horizontal federated learning in practice.
        \item Some solid examples of non-IID data distribution in practical applications.
    \end{itemize}
    
    \item Why is non-IID federated learning a challenging problem? 
    The \textbf{Statistical challenge} brought by non-IID data distribution has the following \textbf{two conflicting aspects}:
    \begin{itemize}
        \item \textbf{Personalization}: When testing, using one model for all prediction tasks on clients with different data distributions usually leads to degenerated prediction performance due to the lack of customization. Using local training or an additional local fine tuning step for customization cannot effectively tackle this problem due to the lack of inter-client collaboration.
        \item \textbf{Collaboration}: When training, not all clients should collaborate with equal strength. Blindly collaborating between clients with equal strength is hard to converge to a good point and, sometimes, even diverges. The theoretical analysis of schemes which perform several local update steps before a communication round is significantly more challenging than those using a single SGD step as in mini-batch SGD. [We use penalty decomposition method to smoothly tackle this problem, such that we can fit typical multi-task learning into a federated learning framework with convergence guarantee.]
        \item \textbf{Dropped Clients and stragglers}
    \end{itemize}
\end{itemize}
}

\nop{
Nevertheless, since local training is possible, it becomes feasible for each client to have a customized model. This approach can turn the non-IID problem from a bug to a feature, almost literally — since each client has its own model, the client’s identity effectively parameterizes the model, rendering some pathological but degenerate non-IID distributions trivial.
}

\nop{
\textbf{Contributions}:
\begin{itemize}
    \item We formulate the task of personalized federated learning on non-IID data as a multi-task learning problem that allows each client to have its own customized model, and encourages clients to adaptively collaborate in clusters by forcing similar client models to become more similar during training.
    
    \item We propose federated attentive message passing (FedAMP) to tackle the multi-task learning problem without infringing the data privacy of any client.
    FedAMP adopts a quadratic penalty optimization approach that binds a shadow model with each client's real model by a proximal term, and conducts inter-client collaboration by attentively passing messages between similar shadow models in a non-linear fashion.
    The non-linear attentive message passing adaptively discovers the underlying collaboration clusters of clients without knowing the number of clusters in prior. This boosts the effectiveness of federated collaboration and leads to the outstanding performance of FedAMP.
    We also proved the convergence of FedAMP for convex client models, and experimentally demonstrated its good practical performance for clients using deep neural networks.
    
    \nop{
    For clients using deep neural networks, we experimentally demonstrate the effectiveness of attentive message passing in achieving high  adaptively discovering the collaboration clusters of clients without knowing the number of clusters in prior.
    }
    
    
    \nop{
    For each client, FedAMP introduces a proxy model that is entangled with the client's real model by a proximal term. The proxy models are used to conduct inter-client collaboration on the cloud.
    }
    \nop{
    \item We propose a personalized federated learning model named Federated Attentive Message Passing (FedAMP) to effectively tackle the non-IID challenge of horizontal federated learning by allowing each client to have its own customized model and attentively passing messages between clients with similar models to make similar client models collaborate more in a non-linear fashion during the training.
    FedAMP adaptively discovers the underlying clustering structure of clients without knowing the number of clusters in prior.
    }
    \nop{The AMP is achieved by ??}
    \nop{
    (non-linearity, ) between the models of clients with different data distributions. \textbf{The key ideas are let each client have its own model, and make similar models collaborate more.}
    }
    \nop{
    \item Our federated learning model is based on a new optimization formulation with nonlinear distance metric. To solve the optimization formulation in setting of federated learning, we introduce a proxy variable and apply a penalty method. Theoretical guarantee. In our algorithm the proxy variable servers as a shadow client model to attentively pass messages between clients in a federated manner. Generate to some simple cases under some situation.
    }
    
    \nop{
     penalty decomposition method to smoothly port general multi-task learning to federated learning framework.}
    
    \item We proposed a Heuristic FedAMP to further improve the practical performance of FedAMP by: a) directly passing messages between real models instead of shadow models to enforce more direct collaboration between clients, and b) measuring the similarity between client models by cosine similarity instead of inner product to avoid scaling problems in the extremely high dimensional space of model parameters.
    
    \nop{
    \item Inspired by the attentive collaboration based on non-linear message passing, and also to tackle the sensitivity to parameters (reason) and (instableness), we developed a simple, efficient and (robust) heuristic method to achieve outstanding practical performance for non-IID horizontal federated learning.}
    
    \nop{
    Explicitly enforce more direct collaboration between clients.
    
    Avoid scaling problem by substituting euclidean distance with cosine similarity to better adapt to high dimensional model parameters.
    }
    
    \nop{
    \item We prove that for simple models, such as linear regression, the attentive collaboration is evaluating the canonical correlation between the data distributions of different clients, and passing stronger messages between clients with more similar data distributions in an non-linear fashion to aggregate their models.
    }
    \nop{
    \item We conduct extensive experiments to demonstrate the effectiveness and efficiency of the proposed methods.
    }
    \nop{
    \item *Naturally handle drop off.
    }
\end{itemize}
}

